\documentclass[11pt]{article}

% ---------- packages ----------
\usepackage[letterpaper,margin=1in]{geometry}
\usepackage{microtype}
\usepackage{parskip}
\usepackage{enumitem}
\setlist{itemsep=2pt,topsep=4pt}
\usepackage[hidelinks]{hyperref}
\usepackage{titlesec}
\usepackage{xcolor}
\usepackage{booktabs}

% ---------- typography ----------
\definecolor{maroon}{HTML}{501214}
\definecolor{gold}{HTML}{D7BD8A}
\definecolor{ink}{HTML}{1A1A1A}

\titleformat{\section}
  {\normalfont\large\bfseries\color{ink}}
  {\thesection.}{0.6em}{}
\titleformat{\subsection}
  {\normalfont\normalsize\bfseries\color{ink}}
  {\thesubsection}{0.6em}{}
\titlespacing*{\section}{0pt}{1.6em}{0.5em}
\titlespacing*{\subsection}{0pt}{1.2em}{0.3em}

% ---------- document ----------
\title{\vspace{-1em}\textbf{Southbound 35}\\
\large A long-term plan for the blog\\[0.3em]
\normalsize \textit{Public finance and economic development on the Texas corridor}}
\author{W.\ Scott Langford\\Texas State University}
\date{June 2026}

\begin{document}
\maketitle
\thispagestyle{empty}

\vspace{-1.5em}
\noindent\textit{This document is an internal planning memo. It sets
the editorial vision, audience, pillar topics, series architecture,
cadence, standards, and growth path for the blog over a three-to-five
year horizon. It is the document I want to be able to point at when
the question ``what is this blog for, exactly?'' comes up --- whether
from a reader, a coauthor, a dean, or me.}

\vspace{1em}
\hrule
\tableofcontents
\hrule

% ============================================================
\section{Vision and positioning}
% ============================================================

Southbound 35 exists to do three things that the existing media
landscape around Texas public finance does not do consistently:

\begin{enumerate}[label=\textbf{(\arabic*)}]
\item Explain the \emph{mechanics} of state and local public finance
in Texas at a level a non-specialist reader can follow, but with
sourcing and replication code that a specialist would not be
embarrassed by.

\item Produce \emph{small-scope analyses} on questions that are too
narrow for an academic paper but too substantive for a Twitter
thread --- the empirical equivalent of a working note, deliberately
addressed to practitioners as well as scholars.

\item Trace the \emph{economic and institutional history of the
I-35 corridor} between Austin and San Antonio, an eighty-mile
strip of land that has become one of the densest growth corridors
in North America while remaining badly under-documented in the
public-facing literature.
\end{enumerate}

The blog's tagline --- \textit{Public finance, economic
development, and applied analysis on the Texas corridor} ---
names the scope deliberately. ``Applied analysis'' is the
umbrella that holds the methodological and program-evaluation
work; ``the Texas corridor'' anchors the empirical setting.
Posts outside these four topics --- sports, religion, off-topic
political commentary --- are out of scope. An earlier editorial
posture allowed ``detours'' into those areas; that posture is
retired as of June 2026.

What Southbound 35 is \emph{not}: an opinion column, a consulting
brochure, a vehicle for academic CV-line-padding, or a substitute
for the academic publications. It is a \emph{public-scholarship
output}, which on the modern academic CV sits in a defined and
defensible category between teaching and research.

% ============================================================
\section{Audience}
% ============================================================

The blog has three concentric audiences, listed from inside to
outside.

\subsection{Inner: practitioners and local officials}

The intended primary reader is someone who works inside or
adjacent to the systems the blog covers --- a school district CFO,
a county auditor, a municipal finance director, a TML or TAC
staffer, a TxDOT or TWDB analyst, a legislative budget board
analyst, a Texas Comptroller staff economist, a Greater Austin
or San Antonio Chamber staffer. These readers know the systems
already; what they value is having the mechanics laid out
cleanly, with the citations they can use in their own briefings,
and the occasional original analysis on a question their
day-to-day workload doesn't give them time to investigate.

This audience is small. It is also disproportionately the audience
that makes consulting requests, suggests new post topics, and
forwards the blog to other practitioners. Almost every quality
decision about the blog should be made with this reader in mind.

\subsection{Middle: engaged residents of the corridor}

The second audience is residents of the I-35 corridor who care
enough about local public finance to read more than 800 words
about it. This includes school board members, city council members,
ISD parents during bond seasons, MUD board members, journalists
covering Texas local government, and the long-tail of Texans who
follow Comptroller releases and TEA reports as a civic hobby.

This audience is much larger than the inner audience and is the
group most likely to share posts on social media, sign up for the
email list, and ask thoughtful questions. The writing voice should
be calibrated for them: assume intelligence and curiosity, do not
assume specialist vocabulary.

\subsection{Outer: scholars and the general public}

The outer ring is the academic community in public administration,
regional economics, public health policy, and Texas studies; plus
the broader public who arrives at the blog through a single shared
post or a news citation. The goal for this audience is not to
convert them to regular readers but to be a credible reference
when they need one.

\subsection{What I am not optimizing for}

\textit{Virality, traffic for its own sake, partisan engagement,
or content velocity.} A blog post that ten engaged practitioners
read and use in their work is worth more than a post that 10{,}000
casual readers click on. The blog's reputation will be made by the
hundred most thoughtful posts I publish over five years, not by the
ten that happen to go briefly viral.

% ============================================================
\section{Editorial pillars}
% ============================================================

Every post fits one of \textbf{four editorial pillars}. Posts
outside these four pillars are not in scope.

\subsection{The four pillars}

\textbf{Pillar 1 --- Public finance.} Explainer posts on how
specific pieces of Texas state and local public finance actually
work, plus original analyses of finance questions where Texas
data permits. The property-tax mechanics, school recapture, and
Truth-in-Taxation explainers are the prototypes; future pillar-1
work covers sales tax, special districts, state-aid formulas, and
bond mechanics. These posts are the load-bearing wall of the
blog and get linked from every subsequent post that uses the
same concepts.

\textbf{Pillar 2 --- Economic development.} Place-based case
studies of growth, infrastructure, employment, and demographics
along the I-35 corridor. The five-post Hays series and the
forthcoming Hays follow-ons, Comal series, and other concentric
build-out (see \S 5) are pillar-2 work. Posts in this pillar are
where the blog's regional identity lives: a specific county, a
specific city, a specific corridor stretch, examined closely.

\textbf{Pillar 3 --- Econometric and statistical analysis.}
Methodological posts that use Texas data to illustrate or apply
a specific method. Difference-in-differences walkthroughs using
a real Texas policy change; regression discontinuity around a
bond-election threshold; synthetic control on a county-level
economic shock; spatial econometrics on corridor outcomes;
replications and re-analyses of published Texas-focused findings.
The audience for these posts is dual --- the methodologically
trained reader gets a clean worked example, and the
practitioner reader gets to see what the method actually does
on a question they care about. The Spring Break Mortality post
is the prototype.

\textbf{Pillar 4 --- Program evaluation.} Evaluations of specific
Texas state and local programs and policies. The Texas Enterprise
Fund, Type A/B Economic Development Corporations, Tax Increment
Reinvestment Zones, the Texas Enterprise Zone program, school
finance reforms (HB 3 compression, SB 2 levy caps), regional
water-plan implementation, and the various incentive vehicles
local governments operate. Posts in this pillar are explicit
about identification strategy, data limitations, and what would
be needed to draw stronger causal conclusions.

The four pillars overlap. A post on HB 3 compression effects is
both \emph{program evaluation} (was the policy effective on its
own terms?) and \emph{econometric analysis} (here is the
identification strategy). A post on Type A/B EDCs is both
\emph{public finance} (how the financing works) and
\emph{program evaluation} (what the audits show). The categories
are organizing principles, not exclusive bins.

\subsection{What the blog is not}

The blog does not publish:

\begin{itemize}
\item Sports analytics, religion and theology, off-topic
political commentary, or other content outside the four pillars
above. (Four legacy posts published under an earlier editorial
scope --- LIV defectors, PGA Tour clutch, Scheffler/Ecclesiastes,
and Spring Break Mortality --- remain live; only the Spring
Break Mortality post fits the new scope, and it is reclassified
as a pillar-3 econometric analysis.)
\item Partisan commentary or candidate endorsements (see \S 9).
\item Direct critique of named individuals' character or
motivations; decisions and policies are critiqued, not the
people who make them (see \S 9).
\item Content that would scoop or puff active academic work
without coauthor review (see \S 6).
\end{itemize}

\subsection{Pillar mix}

There is no fixed pillar ratio. The natural cadence will be
roughly: half of posts in any given quarter are pillar-2
economic-development case studies (because the county-series
arcs are the longest-running content), a quarter are pillar-1
public-finance mechanics and pillar-3 statistical analysis, and
a quarter are pillar-4 program evaluations. The mix can drift
with the seasonal anchors (\S 4) --- e.g.\ pillar-1
Truth-in-Taxation work concentrates in August--September.

% ============================================================
\section{Publishing cadence and the annual rhythm}
% ============================================================

\subsection{Weekly cadence}

The publication slot is \textbf{Monday at 5:00 a.m.\ Central}. This
is locked in by the scheduled rebuild workflow on the website and
the self-hosted mailer that fires a few minutes later. The choice
of slot is deliberate:

\begin{itemize}
\item \textbf{Monday} captures the practitioner workweek-start
attention window. A city manager reading email at 7:30 a.m.\ on
Monday is a better target than a city manager reading at 4:00 p.m.\
on Friday.
\item \textbf{5:00 a.m.\ Central} is late enough that the post is
ready by the start of the workday on the East Coast and early
enough not to be buried by Monday-morning newsletters.
\end{itemize}

The cadence is \emph{one post per week}, every week. Skipping weeks
is much more costly than publishing a less-than-perfect post on
time. The hardest discipline is preferring a B+ post that ships
Monday to an A post that ships ``some Friday.''

\subsection{Annual rhythm}

The Texas public finance calendar has a strong seasonal rhythm,
and the blog should ride it:

\begin{itemize}
\item \textbf{January--May (odd years).} Legislative session. Daily
news flow is heavy; the blog's contribution is structural
explainers (what each finance-relevant bill actually does) rather
than play-by-play.

\item \textbf{June--July.} Comptroller's biennial revenue estimate
and the certified estimate refresh. Sales tax monthly data start
showing the post-spring picture. Good window for state-revenue
posts.

\item \textbf{August--September.} School district budget adoption
and tax rate setting. Truth-in-Taxation hearings statewide. The
single best season of the year for explainer posts on local property
taxation.

\item \textbf{October--November.} Bond elections (uniform election
date is the first Tuesday after the first Monday in November). One
post per cycle on the largest corridor bond elections, with results
follow-ups.

\item \textbf{December.} Census ACS five-year estimates released.
Annual ``year in review'' post drafted in the last week of December.
\end{itemize}

A good content calendar is not a list of fifty-two specific topics
locked a year in advance. It is a list of the seasonally-anchored
slots (about fifteen per year) plus a backlog of pillar topics that
can be slotted into the open weeks.

% ============================================================
\section{Series architecture}
% ============================================================

The blog's distinctive intellectual structure --- the thing that
separates it from a column or a newsletter --- is the
\textbf{multi-post series}. Each series follows a standard arc that
lets a reader build a coherent picture of a place or a topic over
four to six posts.

\subsection{The county series template}

The Hays County series established the template. Five posts:

\begin{enumerate}
\item \textbf{Growth.} The demographic and economic story. Why has
this county been growing? Where are the new residents coming from?
What is the labor-market base?
\item \textbf{Projections.} What do credible forecasts say about the
next ten and twenty-five years? Where do they disagree, and on what?
\item \textbf{Water.} The single binding resource constraint in
most fast-growing Texas counties. Sources, providers, regulatory
geography, and what the next twenty years require.
\item \textbf{Schools.} The largest line on the property tax bill
and the most politically charged piece of local governance.
Enrollment, finance, recapture status, bond capacity.
\item \textbf{Governance.} Who actually governs this place? County,
cities, ISDs, special districts, ETJs. The institutional plumbing.
\end{enumerate}

The order in which the template is applied to additional
counties follows a deliberate \emph{geographic-proximity} rule:
the blog is written from San Marcos, Texas State University sits
in Hays County, and the public-scholarship value of the work is
highest where the author can ground-truth the institutional
landscape. Series are therefore commissioned in concentric rings
out from San Marcos:

\textbf{Ring 0 --- Hays itself.} Deepen the existing five-post
series with follow-on case studies (San Marcos and the
Texas State enrollment-finance feedback; the Dripping Springs /
western-Hays Hill-Country edge; Kyle/Buda as the corridor's
fastest-growing pair; Hays CISD vs. San Marcos CISD vs.
Dripping Springs ISD as a three-district school-finance
comparison).

\textbf{Ring 1 --- counties directly adjacent to Hays.} In
priority order: Comal (south; same I-35 corridor, Edwards
Aquifer politics, New Braunfels anchor); Caldwell (east; Lockhart
anchor, slower growth, distinct rural-to-exurban transition);
Guadalupe (south/southeast; Seguin and Schertz/Cibolo, Toyota
supplier ecosystem); the southern edge of Travis County
(Dripping Springs / Manchaca / southwest Austin --- the part of
Travis that interacts with Hays daily); and the eastern edge of
Blanco County (the Hill-Country boundary with western Hays).

\textbf{Ring 2 --- the rest of the corridor.} Williamson (the
closest analogue to Hays in growth dynamics, but separated by
two counties); the Bexar Northside (a different governance
pattern --- a single large city extending into a fragmented
exurban edge); and the Bastrop / eastern-Caldwell edge (the
most under-studied piece of the corridor).

\subsection{The pillar series template}

A pillar series is a sequence of mechanics posts on a single
finance system, building up the reader's vocabulary cumulatively.
The currently-drafted three-post sequence on property taxation
(mechanics, recapture, Truth-in-Taxation) is the first. Future
pillar series will cover, in order:

\begin{itemize}
\item Sales tax mechanics: state, local-option, Type A/B
economic development corporations, MTA/ATD slices.
\item Special districts: MUDs, ESDs, hospital districts,
groundwater conservation districts. The least-understood layer of
Texas government.
\item State-aid formulas: school FSP, county road \& bridge
distributions, the various sales-tax holdback mechanisms.
\item Bond mechanics: voter-approval, debt service, taxable
versus tax-exempt issuance, the legal opinions and what they
mean.
\end{itemize}

% ============================================================
\section{Quality standards and sourcing}
% ============================================================

The blog's credibility is its most valuable asset. Six rules:

\begin{enumerate}
\item \textbf{Every chart cites its source in the chart footer.}
Not in a paragraph below, not in a parenthetical, not in the
replication code. In the chart itself.

\item \textbf{Every post that uses non-trivial data points to a
replication package} under \texttt{southbound-35/posts/}. The
package is committed to a public repo on the same calendar day as
publication.

\item \textbf{Illustrative figures are labeled as such.} When a
chart uses approximate values or transcribed summary statistics
because a clean machine-readable source is unavailable, the
chart's source line says so.

\item \textbf{Claims about specific districts, officials, or
firms are checked against named primary sources.} TEA, Comptroller,
Census, BEA, FDIC, the agency's own adopted budget. Press
coverage of a finance event is a starting point, not a citation.

\item \textbf{Errors are corrected in-line with a dated note.}
Not silently. The note says what was wrong, when it was changed,
and (when relevant) who flagged it.

\item \textbf{Posts that touch active academic work go through the
same coauthor review as the academic work.} If a paper is in
review at a journal and the blog post repeats or anticipates its
finding, the coauthors see the post before it publishes.
\end{enumerate}

% ============================================================
\section{Restructuring the existing posts}
% ============================================================

A blog written week by week without a planning document
accumulates inconsistencies in its metadata. The first six months
of Southbound 35 are no exception. The restructuring proposed
here is metadata-only --- it does not change post content or
visual rendering --- and is intended to make the next two years
of content navigable, both for readers and for the site's own
indexing and series widgets.

\subsection{Front-matter conventions}

Every post should declare four pieces of front matter, beyond
title and date:

\begin{itemize}
\item \texttt{categories} --- one of four pillar buckets:
\texttt{public-finance}, \texttt{economic-development},
\texttt{econometric-analysis}, or \texttt{program-evaluation}.
(A post may carry more than one category, but not more than two
--- many posts will naturally land on a pillar boundary,
e.g.\ a HB 3 compression piece is both
\texttt{public-finance} and \texttt{program-evaluation}.) The
legacy \texttt{detours} value remains on the four pre-June-2026
posts that fall outside the new scope; it is not used for new
content.

\item \texttt{series} --- a slug identifying any multi-post series
the post belongs to (\texttt{hays-county},
\texttt{property-tax-mechanics}, \texttt{i35-corridor},
\texttt{special-districts}, etc.).
Posts that do not belong to a series omit the field.

\item \texttt{kind} --- one of: \texttt{explainer},
\texttt{original-analysis}, \texttt{case-study}, or
\texttt{evaluation}. The legacy \texttt{commentary} and
\texttt{detour} values are not used for new content.

\item \texttt{tags} --- free-form topic labels (Texas, water,
schools, golf, theology, etc.).
\end{itemize}

The point of separating \texttt{categories}, \texttt{series},
\texttt{kind}, and \texttt{tags} is to make different navigation
widgets cheap to add later: a series-spine widget that shows
``you are reading post 3 of 5 in the Hays County series'';
a kind-based filter on the blog landing page; a
related-by-tag panel at post foot.

\subsection{Existing posts: corrected metadata}

The following table shows the proposed metadata for the nine
existing live posts. Changes are minor and reversible.

\begin{center}
\small
\begin{tabular}{p{1.6in}p{0.9in}p{1.0in}p{1.0in}}
\toprule
Title & Categories & Series & Kind \\
\midrule
Hays Growth        & econ-dev, pub-fin & hays-county    & case-study \\
Hays Projections   & econ-dev, pub-fin & hays-county    & case-study \\
Hays Water         & econ-dev, pub-fin & hays-county    & case-study \\
Hays Schools       & econ-dev, pub-fin & hays-county    & case-study \\
Hays Governance    & econ-dev, pub-fin & hays-county    & case-study \\
\midrule
Spring Break Mortality & detours       & ---            & original-analysis \\
LIV Defectors          & detours       & ---            & original-analysis \\
Golf Clutch            & detours       & ---            & original-analysis \\
Scheffler / Ecclesiastes & detours     & ---            & commentary \\
\bottomrule
\end{tabular}
\end{center}

\subsection{Series-aware permalinks}

Permalinks should not be changed retroactively (that breaks
backlinks). Going forward, however, series posts should use a
permalink template like
\texttt{/series/hays-county/water/} rather than
\texttt{/posts/2026/05/hays-county-water/}, with redirects from
the dated path. This makes the series structure visible in URLs.

% ============================================================
\section{Distribution, audience-building, and email}
% ============================================================

The blog has three distribution channels: the website itself, an
email list, and social-media reposting.

\subsection{Email is the only one that compounds}

Of the three, only the email list is owned and durable. Twitter /
X audiences are subject to platform decisions outside the blog's
control; Bluesky is more aligned but still small; LinkedIn is
useful for practitioner distribution but the algorithm rewards
content patterns the blog should not adopt.

The self-hosted mailer (companion repo \texttt{southbound-35},
directory \texttt{mailer/}) sends each subscriber an individual
email Monday morning. No BCC, no shared header, no third-party
provider. Subscribers are stored in a flat file on the desktop
machine. The list is not shared, sold, or analyzed beyond the
delivery confirmations.

\textbf{Year-one goal: 200 subscribers, all opt-in.} Two-thirds
of these will come from practitioner networks (TML, TAC, GFOA,
ASPA, NASPAA, the Texas Public Policy Foundation and Center for
Public Policy Priorities, the Texas state agencies). One-third
will come from the broader social-media surface.

\textbf{Year-three goal: 1{,}000 subscribers}, with a stable
weekly open rate above 50\%. At 1{,}000 practitioner-heavy
subscribers, the blog has reached its natural audience ceiling
for Texas-only content. Further growth would require either
national-scope content (something I do not currently plan to
write) or a successor product.

\subsection{Social channels}

The blog's primary social-media presence is Bluesky
(\texttt{@scottlangford2.bsky.social}) for academic-adjacent
distribution and LinkedIn for practitioner reach. Twitter / X is
posted to but not relied on. The accounts repost the blog's
weekly post once, with the post's pull quote and the headline
chart. They do not engage in unrelated commentary, threads, or
political conversation under the same handle. This is a
deliberate restraint.

\subsection{Cross-posting and syndication}

The blog will not be syndicated to Substack, Medium, or other
third-party platforms. The reason is partly editorial control
and partly the principle that the canonical version is at the
canonical URL.

It will, however, be excerpted and reprinted with permission by
Texas-focused outlets that want to. The first such relationship
to develop is with the Texas State University Hobby School (or
analogue), the LBJ School at UT, and the SA Express-News
Editorial Board, all of which have institutional interests in
publishing finance-explainer pieces aimed at engaged residents.

% ============================================================
\section{Risks, boundaries, and ethics}
% ============================================================

The blog has several risk surfaces. None is fatal; all need to
be managed deliberately.

\subsection{Conflict of interest with consulting}

I do small-scope consulting work for local governments and
development organizations. The risk is that a blog post praises
or criticizes a client in a way that creates the appearance of
quid pro quo. The rule: \textbf{no post on a current consulting
client's jurisdiction or work product during the engagement, and
for six months after.} Disclose past engagements in any post that
discusses the relevant jurisdiction.

\subsection{Conflict of interest with research}

The blog should not be used to scoop my own academic papers. It
also should not be used to puff them. The rule: \textbf{academic
work appears on the blog only after the underlying paper is in
working-paper form or beyond}, and a blog post about it is
explicit that it is a public-facing summary.

\subsection{Defamation and named-individual claims}

The blog will name elected officials, civil servants, and public
figures in their official capacities, and will critique
\emph{decisions} freely. It will not critique individuals'
character or motivations. The rule: \textbf{disagree with the
decision, not the person who made it.}

\subsection{Political affiliation}

The blog does not endorse candidates or parties and does not
publish commentary that could be mistaken for endorsement. This
is a professional-norm restriction, not a personal one.

\subsection{Texas State University and academic position}

The site footer makes clear that the blog is a personal
publication. University communications are not consulted on
content. The university's name and brand are not used as
endorsement. Standard academic-freedom protections apply, and so
do the standard academic-norm restrictions on partisan activity.

% ============================================================
\section{Measurement and evaluation}
% ============================================================

The blog is measured on four axes, listed in order of importance.

\subsection{Substance reach (most important)}

Are practitioners citing or forwarding the blog? Concretely:
\begin{itemize}
\item Citations of specific posts in TEA, Comptroller, TML, TAC
publications.
\item Inclusions in regional planning documents.
\item Practitioner referrals (someone says ``I came here from
[a colleague's] forward'').
\item Press coverage that links to the blog by name.
\end{itemize}

This is the metric I actually care about. It is hard to
instrument, but the qualitative signal is clear when present.

\subsection{Email list (next most important)}

Subscriber count, weekly open rate, click-through to source
material. The mailer logs these.

\subsection{Site traffic (least important)}

Plausible Analytics provides cookieless aggregate counts. The
metric is monitored for diagnostic purposes (a 90\% week-over-week
drop indicates a broken deploy; a 10x spike indicates the post
was shared somewhere) but is not optimized for.

\subsection{Personal: am I learning?}

The deepest justification for the blog, beyond reach and audience,
is that writing it teaches me things about Texas public finance
that I would not otherwise know. The honest measure is whether,
six months from now, I can answer questions about MUD finance,
recapture, or corridor traffic that I could not have answered
today. If yes, the blog is doing its work even if the audience
metrics are flat.

% ============================================================
\section{Long-horizon outputs}
% ============================================================

\subsection{Year 1 (June 2026 -- May 2027)}

Steady weekly cadence. Two completed pillar series
(property-tax mechanics, special districts). Two completed
corridor series: a deeper Hays follow-on track (San Marcos /
TXST finance feedback; the western Hays Hill-Country edge), and
the Comal series. Reach 200 email subscribers.
First practitioner citation by a state agency. First press
mention by a corridor newspaper.

\subsection{Year 2 (June 2027 -- May 2028)}

One pillar series and one corridor series per quarter. Reach
500 email subscribers. First non-Texas reader base (regional
economists, scholars of state-and-local finance more broadly).
Begin compiling the explainer posts into a structured
reference, with the working title \emph{Texas Public Finance:
A Practitioner's Guide}, distributed as a free PDF.

\subsection{Year 3--5 (June 2028 onward)}

Reach 1{,}000 email subscribers and approach the natural ceiling
for Texas-only content. Publish the practitioner's guide as a
formal open-access book through Texas State University Press or
an analogue. Begin a second blog product (open question whether
this is national in scope or a separate Texas-focused publication
with a different editorial voice).

\textbf{Five-year vision in one sentence:} \emph{Southbound 35 is
the canonical public-facing reference on the institutional plumbing
of fast-growing Texas, written from inside the academy, for the
practitioners who run the places.}

% ============================================================
\section{Implementation checklist}
% ============================================================

The concrete near-term to-do list, in priority order:

\begin{enumerate}
\item Normalize front matter on the nine existing posts
(\texttt{series}, \texttt{kind} fields) per the
table in \S 7.2. This is the immediate next action and is
metadata-only.

\item Draft the three queued posts to publication-ready and review
them on iPad before pushing live:
\begin{itemize}
\item Property tax mechanics 101
\item Eanes ISD recapture
\item I-35 Austin--SA corridor history
\end{itemize}

\item Add a series-spine widget to the post layout that
renders ``Post 3 of 5: Hays County series'' when \texttt{series}
is present. This is theme work and should be scoped before
implementation.

\item Draft the next Hays follow-on posts (San Marcos and the
Texas State enrollment-finance feedback; Kyle/Buda as a paired
case study; a three-ISD school-finance comparison across Hays
CISD, San Marcos CISD, and Dripping Springs ISD), and the first
two Comal County series posts. Review each PDF before the
cadence date.

\item Build a manual content-calendar tracker (a small JSON file
under \texttt{plan/} listing the next twenty Mondays and their
slotted topic) and review it monthly.

\item Develop a one-page editorial policy document derived from
\S 6 and \S 9 of this plan, suitable for sharing with collaborators
and coauthors.
\end{enumerate}

% ============================================================
\appendix
\section{Appendix: front-matter template}
% ============================================================

The canonical front-matter for new posts is:

\begin{verbatim}
---
title: 'Post title in title case'
date: YYYY-MM-DD
permalink: /posts/YYYY/MM/slug/   # legacy form; series posts
                                  # also get /series/slug/post-slug/
related: false
categories:
  - public-finance               # one of four pillars:
                                 # public-finance, economic-development,
                                 # econometric-analysis, program-evaluation
series: hays-county              # slug; omit if no series
kind: case-study                 # explainer | original-analysis |
                                 # case-study | evaluation
tags:
  - Texas
  - Hays County
  - public finance
excerpt: 'One- or two-sentence summary suitable for use as
  the post deck in print and the social-share description.'
---
\end{verbatim}

\end{document}
